\subsection{Light-Emitting Diode (LED)}
\markright{Light-Emitting Diode (LED)}

\subsubsection*{Definition}
A light-emitting diode (LED) is a forward-biased p--n junction that converts electrical
energy directly into light through a process called \textit{electroluminescence}
\cite{sedra2020}. When electrons and holes recombine across the junction, they release
energy in the form of photons. The wavelength---and thus color---of the emitted light
depends on the semiconductor material's bandgap energy. Red and infrared LEDs use
AlGaAs or GaAsP; green and yellow use GaP or AlInGaP; high-brightness blue and white
LEDs (commercially available since the 1990s, earning Akasaki, Amano, and Nakamura the
2014 Nobel Prize in Physics) use InGaN \cite{razavi2013}.

\labimage{lab-01/assets/led.jpeg}%
  {Standard 5\,mm through-hole LEDs in red, yellow, and green: the longer
   lead is the anode; the flat edge on the base identifies the cathode.}{led_lab}

\subsubsection*{Specifications}
Typical 5\,mm indicator LEDs have: $V_F = 1.8$--$2.2$\,V (red/yellow) or
$3.0$--$3.5$\,V (blue/white), maximum $I_F = 20$\,mA, luminous intensity of
10--1000\,mcd, and viewing angle of 15$^\circ$--60$^\circ$ \cite{horowitz2015}.
High-power LEDs rated at 350\,mA to several amperes are used in lighting applications.

\subsubsection*{Purpose of Use}
In the laboratory, LEDs are used as power and signal indicator lights. They must always
be paired with a current-limiting resistor calculated as
$R = (V_\text{supply} - V_F) / I_F$.

\subsubsection*{Applications}
LEDs appear in indicator panels, numeric and alphanumeric displays, LCD backlighting,
automotive headlamps and tail lights, traffic signals, optical fiber transmitters,
infrared remote controls, and solid-state general lighting \cite{razavi2013}.
